\section{Nghiên cứu con người}

\subsection{Các Phát hiện Human Evaluation Chính}

\begin{table}[H]
\centering
\caption{Kết quả Human Evaluation}
\label{tab:human-studies}
\begin{tabular}{lll}
\toprule
\textbf{Study} & \textbf{N} & \textbf{Phát hiện} \\
\midrule
PersonaEval (Zhou et al., 2025) & -- & Human: 90.8\% speaker-ID vs LLM: $\sim$69\% (gap -21.8pp) \\
RMTBench (2025) & -- & Human annotator agreement $\kappa=0.77$--$0.84$ (16--23\% disagreement) \\
Bickmore \& Picard (2005) & 52 & 22\% churn trong 2 tuần đầu, retention improves to 91\% by week 12 \\
Yang \& Oshio (2025) & 200 & Familiarity +24\%, trust +19\% over 4 weeks \\
\bottomrule
\end{tabular}
\end{table}

\subsection{Các Mẫu Relationship Human-AI}

\begin{enumerate}
    \item \textbf{Mô hình Ưu tiên Quen thuộc}: Mức quen thuộc tăng nhanh hơn trust
    \item \textbf{Kiểu gắn bó quan trọng}: Anxious attachment $\to$ intimacy mạnh hơn; Avoidant $\to$ ảnh hưởng tiêu cực
    \item \textbf{Tác động lỗi}: Một lỗi nghiêm trọng có thể giảm trust 23--37\%; sửa chữa giảm thiệt hại $\sim$40\%
    \item \textbf{Mô hình phục hồi}: Hàm mũ, không tuyến tính
\end{enumerate}

% =============================================================================
% 22. LONGITUDINAL STUDIES